% !TEX root = ../main.tex
\clearpage
\section{创新点与竞争分析}
\subsection*{金融研究方法上的改进}\label{page:33}
\ifdraftversion{\noindent\small\color{gray}写作页位 33\quad 负责人 武桐宇\par}\fi

\textbf{本页结论：}项目的方法改进不在于宣称发现新的金融规律，而在于把低空研究中易断裂的证据、产业关系和评分规则组织成可检查、可撤回、可复算的工作流。

方法的起点是“先给来源，再给结论”。一则事件若不能定位原文、适用对象和日期，便不进入正式集；一条公司关系若未核验，便不参与暴露度；一组财务字段若无统一口径，便不计算财务重要性代理。这样做牺牲了覆盖面，却把每一步的责任边界暴露给复核者。

\begin{table}[htbp]
\centering
\caption{研究方法设计与可验证问题的对应}
\label{tab:wty-method-improvement}
\small
\begin{tabularx}{\linewidth}{p{0.22\linewidth}X p{0.25\linewidth}}
\toprule
设计 & 解决的具体问题 & 可验证依据 \\
\midrule
事件证据卡 & 新闻转述难回到原文与日期 & 链接、位置、状态与来源账本 \\
产业关系图谱 & 产业相邻被误写为公司关联 & 每条边和映射的来源、状态与撤回记录 \\
公式化强度评分 & 结论难以比较或复算 & 固定 $E,S,F$ 输入、方向字段与手算例 \\
版本化导出 & 结果更新后无法解释差异 & 数据包、规则、生成日期和差异日志 \\
\bottomrule
\end{tabularx}
\end{table}

该方法仍需由真实、授权数据和人工核验补全。它不声称已经完成因果识别、收益预测或大模型自动研究，而是为这些后续工作提供可审查的数据与规则基础。
